\chapter{词法分析}


\section{词法分析概述}


词法分析(morphological analysis)是 NLP 系统接收到输入字符串之后, 在进入句子层面的句法分析和语义分析之前, 对输入字符串进行的词一级处理.

可以把它放在一个基本处理流水线中理解:

\begin{Verbatim}[breaklines=true,breakanywhere=true,fontsize=\small]
字符串 -> 词法分析 -> 句法分析 -> 语义分析
\end{Verbatim}

也就是说, 词法分析关心的是:
系统如何从连续字符串中识别词、词的边界、词的类别以及文本中的特殊实体.

课件中列出的主要任务包括:

\begin{itemize}
\item 汉语自动分词.
\item 英语形态还原.
\item 命名实体识别.
\item 词性标注.
\end{itemize}

其中汉语 NLP 尤其依赖词法分析.
因为汉语书面文本通常没有天然空格分隔词语, 如果分词结果错误, 后续的句法分析、语义分析、检索、分类和问答都会受到影响.

\section{汉语自动分词}


\subsection{基本概念}


汉语自动分词是让计算机系统在汉语文本中的词与词之间自动加上边界标记.

例如:

\begin{Verbatim}[breaklines=true,breakanywhere=true,fontsize=\small]
我们在上自然语言处理课程。
\end{Verbatim}

可以切分为:

\begin{Verbatim}[breaklines=true,breakanywhere=true,fontsize=\small]
我们 / 在 / 上 / 自然语言处理 / 课程 / 。
\end{Verbatim}

分词看似只是加边界, 实际上是在回答“什么应该被当作一个词”这个问题.
如果边界划分不同, 句子的结构和含义也可能发生变化.

\subsection{主要困难}


汉语自动分词的困难主要来自三个方面:

\begin{itemize}
\item 分词规范不明确.
\item 切分歧义普遍.
\item 未登录词频繁出现.
\end{itemize}

\subsubsection{分词规范不明确}


分词规范用于明确什么应该被视作独立单元切分出来.
一般认为汉语中应被切分出的独立单元包括两类:

\begin{itemize}
\item 有明确内涵定义的词语.
\item 紧密、稳定、内涵独立的词组.
\end{itemize}

困难在于, “紧密”“稳定”“内涵独立”并不总是容易判断.
例如“大白菜”“小媳妇”“白花”等表达, 在不同任务或不同词表规范下可能有不同切法.

因此, 汉语分词并不存在一个完全唯一、天然正确的答案.
一个分词系统的结果往往依赖其采用的词表、标注规范和应用目标.

\subsubsection{切分歧义}


切分歧义是指同一汉字串存在多种从形态层面看都合理的切分方式.

一种常见类型是交集型歧义.
如果汉字串 \texttt{AJB} 中, \texttt{AJ} 和 \texttt{JB} 都是词, 则 \texttt{AJB} 可形成交集型歧义.

例如:

\begin{Verbatim}[breaklines=true,breakanywhere=true,fontsize=\small]
乒乓球拍卖完了
\end{Verbatim}

可以理解为:

\begin{Verbatim}[breaklines=true,breakanywhere=true,fontsize=\small]
乒乓球拍 / 卖 / 完 / 了
乒乓球 / 拍卖 / 完 / 了
\end{Verbatim}

另一种常见类型是组合型歧义.
如果汉字串 \texttt{AB} 中, \texttt{A}、\texttt{B}、\texttt{AB} 都是词, 则 \texttt{AB} 可形成组合型歧义.

例如“马上”:

\begin{Verbatim}[breaklines=true,breakanywhere=true,fontsize=\small]
他 / 从 / 马 / 上 / 跳 / 下来
我 / 马上 / 就 / 来
\end{Verbatim}

第一句中“马”和“上”应分开, 第二句中“马上”应整体作为副词.
因此, 分词需要利用上下文, 不能只看局部字串是否在词典中.

\subsubsection{未登录词}


未登录词是指不在词表中的词.
它们会频繁出现, 尤其包括:

\begin{itemize}
\item 新词, 如“博客”“房奴”“内卷”等.
\item 专有名词, 如人名、地名、机构名、时间、数字等.
\item 专业名词, 如医学、金融、工程等领域中的术语.
\end{itemize}

未登录词对分词影响很大.
如果系统完全依赖固定词典, 它往往会把新词错误拆开, 或把实体边界切错.
这也是统计模型、序列标注模型和深度学习模型后来成为主流的重要原因.

\subsection{最大匹配法}


最大匹配法是一种基于词典的机械切分方法.
其基本思想是: 在当前位置尽量匹配词典中最长的词, 匹配成功就切出该词, 否则把当前字作为单字词处理.

按照扫描方向可分为:

\begin{itemize}
\item 正向最大匹配: 从左到右扫描.
\item 反向最大匹配: 从右到左扫描.
\end{itemize}

正向最大匹配可以概括为:

\begin{enumerate}
\item[1.] 从当前位置开始取不超过最大词长的字串.
\item[2.] 若该字串在词典中, 则切出该词.
\item[3.] 若不在词典中, 则缩短一个字继续尝试.
\item[4.] 若仍无匹配, 则把当前字作为单字词.
\item[5.] 移动到下一个未处理位置并重复.
\end{enumerate}

最大匹配法实现简单, 速度快, 在词典质量较高且文本规范的场景中有一定效果.
但它有两个明显问题:

\begin{itemize}
\item 不能识别未登录词.
\item 总是追求局部最长匹配, 可能错过全局最优切分.
\end{itemize}

因此, 最大匹配法适合作为简单基线, 但仅凭词典和局部匹配, 很难稳定解决歧义和未登录成分.

\subsection{最短路径法}


最短路径法也是基于词典的切分方法, 但它不只在当前位置做局部最长选择, 而是把所有可能切分组织成一张词图.

基本思路是:

\begin{itemize}
\item 把句子中每个字间位置看作图中的节点.
\item 若某个连续字串在词典中, 就在对应起止位置之间连一条边.
\item 每条边表示一种可能词项.
\item 从句首到句尾的一条路径对应一种完整切分.
\end{itemize}

最短路径法认为: 词语数量最少的切分方式就是较优结果.
这相当于在词图中寻找边数最少的路径.

例如课件中的输入字串:

\begin{Verbatim}[breaklines=true,breakanywhere=true,fontsize=\small]
他只会诊断一般的疾病。
\end{Verbatim}

若词典中同时存在\texttt{只会}、\texttt{诊断}、\texttt{会诊}、\texttt{一般}、\texttt{疾病}等词,
系统可以得到多条候选路径:

\begin{Verbatim}[breaklines=true,breakanywhere=true,fontsize=\small]
他 / 只会 / 诊断 / 一般 / 的 / 疾病 / 。   词个数: 7
他 / 只 / 会诊 / 断 / 一般 / 的 / 疾病 / 。 词个数: 8
\end{Verbatim}

最短路径法会选择词个数更少的第一种切分:

\begin{Verbatim}[breaklines=true,breakanywhere=true,fontsize=\small]
他 / 只会 / 诊断 / 一般 / 的 / 疾病 / 。
\end{Verbatim}

如果多条路径边数相同, 通常还需要结合词频、词长、概率或其他规则进一步消歧.

相比最大匹配法, 最短路径法考虑了句子级别的整体切分, 能避免一部分局部贪心错误.
但它仍依赖词典, 因而同样难以识别未登录词.
如果词图中没有正确词项, 再好的路径搜索也无法恢复它.

\subsection{最大概率法}


最大概率法把分词看作概率最大化问题.
给定原始字符串$s$, 设一种候选切分为:
\[
w = (w_1, w_2, \dots, w_m)
\]

目标是在所有可能切分$\Omega(s)$中选择概率最大的切分:
\[
w^* = \arg\max_{w \in \Omega(s)} P(w | s)
\]

因为原始字符串$s$由切分结果$w$连接而成, 常把问题转化为最大化切分序列本身的概率:
\[
w^* = \arg\max_{w \in \Omega(s)} P(w)
\]

若采用 n-gram 语言模型, 可以近似为:
\[
P(w) \approx \prod_{i=1}^m P(w_i | w_{i-1})
\]

于是, 分词问题就变成:
在词图中寻找一条总概率最高的路径.

最大概率法的优势是能够利用语料统计信息处理歧义.
例如某种切分中的相邻词搭配在语料中更常见, 它就会获得更高概率.
但如果候选词本身没有进入词图, 未登录词问题仍然存在.

\subsection{字分类法}


字分类法把分词转化为序列标注问题.
其基本思想是: 不直接枚举词, 而是从头到尾为句子中的每个字预测一个词位标签.

常见四分类标记为:

\begin{itemize}
\item \texttt{B}: 词的开始字(begin).
\item \texttt{M}: 词的中间字(middle).
\item \texttt{E}: 词的结束字(end).
\item \texttt{S}: 单字词(single).
\end{itemize}

例如分词结果:

\begin{Verbatim}[breaklines=true,breakanywhere=true,fontsize=\small]
乒乓球 / 拍 / 完 / 了
\end{Verbatim}

可以对应为:

\begin{Verbatim}[breaklines=true,breakanywhere=true,fontsize=\small]
乒/B 乓/M 球/E 拍/S 完/S 了/S
\end{Verbatim}

完成逐字标注之后, 再根据标签恢复词边界.
因此, 分词被统一成一个标准序列标注任务.

字分类法通常使用上下文特征.
对当前待分类字$C_0$, 常用模板包括:

\begin{itemize}
\item 单字上下文: $C_k$, 其中$k=-2,-1,0,1,2$.
\item 相邻字组合: $C_k C_{k+1}$, 其中$k=-2,-1,0,1$.
\item 跨字组合: $C_{-1} C_1$.
\item 字符类别特征: 数字、标点、英文字母等.
\end{itemize}

常用模型包括:

\begin{itemize}
\item HMM.
\item CRF.
\item SVM.
\item 神经网络模型.
\end{itemize}

字分类法比纯词典方法更适合处理未登录词.
只要模型学到了构词模式和上下文线索, 即使某个词从未出现在词表中, 也有可能预测出合理边界.

\subsection{当前问题}


汉语分词已经能在规范文本和同领域测试集上达到较高效果, 但仍有明显局限:

\begin{itemize}
\item 跨领域应用时性能不理想, 例如用新闻语料训练的模型处理医学文本.
\item 面对非规范文本时性能下降, 例如微博、短信、网络评论等.
\item 新词、缩写、口语表达和专有名词不断出现, 需要模型具备持续适应能力.
\end{itemize}

因此, 分词不是一次性解决的问题.
实际系统通常需要结合通用模型、领域词典、实体识别和人工反馈来持续维护.

\section{命名实体识别}


\subsection{基本概念}


命名实体识别(named entity recognition, NER)是识别文本中对实体的命名性指称.
具体来说, 它要识别人名、地名、机构名、时间名等专有名词.

例如:

\begin{Verbatim}[breaklines=true,breakanywhere=true,fontsize=\small]
中国乒乓球男队总教练刘国梁出席了备战会议，他指出了当前备战工作的重点。
\end{Verbatim}

其中实体可以是:

\begin{Verbatim}[breaklines=true,breakanywhere=true,fontsize=\small]
刘国梁
\end{Verbatim}

而文本中对同一实体可能有多种指称方式:

\begin{itemize}
\item “中国乒乓球男队总教练”: 名词性指称.
\item “刘国梁”: 命名性指称.
\item “他”: 代词性指称.
\end{itemize}

NER 关注的是命名性指称, 也就是文本中明确以名称形式出现的实体.

\subsection{主要方法}


命名实体识别的主要方法可分为三类:

\begin{itemize}
\item 基于规则的方法.
\item 基于统计的方法.
\item 基于深度学习的方法.
\end{itemize}

\subsubsection{基于规则的方法}


基于规则的方法依赖人工整理的词表和模式.
以人名识别为例, 可以建立姓氏用字规则和上下文识别规则.

姓氏用字可以分为:

\begin{itemize}
\item 几乎总是用作姓氏的字, 如陈、丁、邓、蒋等.
\item 偶然用作姓氏的字, 如曾、都、向、于等.
\item 作姓氏与非姓氏机会接近的字, 如马、黄、张等.
\item 待定字, 如任、房、方等.
\end{itemize}

上下文规则则利用周围词语判断一个字是否更可能作为姓名的一部分.
例如:

\begin{itemize}
\item “老/小 + 姓氏用字”通常可识别人称, 如“小王”“老李”.
\item “姓氏用字 + 工/总”通常可识别人称, 如“张工”“陈总”.
\item “数词 + 可作量词的姓氏用字”通常不识别人名, 如“一周”“七章”.
\end{itemize}

规则方法的优点是可解释、易于注入领域知识.
缺点是维护成本高, 对新场景和复杂上下文的覆盖有限.

\subsubsection{基于统计的方法}


基于统计的方法把 NER 看作字序列标注问题.
以人名识别为例, 可以把每个字标注为:

\begin{itemize}
\item \texttt{PNB}: 人名的起始用字.
\item \texttt{PNI}: 人名的内部用字.
\item \texttt{OUT}: 不属于任何人名实体.
\end{itemize}

例如:

\begin{Verbatim}[breaklines=true,breakanywhere=true,fontsize=\small]
小王积极参加班级的各项活动，因此很受班主任张老师的赏识。
\end{Verbatim}

可以标注为:

\begin{Verbatim}[breaklines=true,breakanywhere=true,fontsize=\small]
小/PNB 王/PNI 积/OUT 极/OUT ... 张/PNB 老/PNI 师/PNI ...
\end{Verbatim}

常用特征包括:

\begin{itemize}
\item 当前字及其左右窗口内的字.
\item 当前字是否为常见姓氏用字.
\item 当前字是否为指界词, 如称谓、动词、标点等.
\item 字符类型特征.
\end{itemize}

常用模型包括 CRF 等序列标注模型.
这类方法比纯规则方法更容易从标注语料中学习上下文模式.

\subsubsection{基于深度学习的方法}


深度学习方法通常不再手工设计大量离散规则, 而是通过神经网络自动学习表示.

典型流程是:

\begin{enumerate}
\item[1.] 将句子中的字或词映射为 embedding.
\item[2.] 将 embedding 序列输入神经网络.
\item[3.] 由神经网络提取上下文特征.
\item[4.] 通过分类层或 CRF 层预测每个位置的实体标签.
\end{enumerate}

常见模型包括:

\begin{itemize}
\item RNN.
\item LSTM.
\item BiLSTM.
\item BERT 等预训练模型.
\end{itemize}

现代 NER 系统常采用 \texttt{B-I-O} 或 \texttt{B-M-E-S} 标注体系.
例如 \texttt{B-PER} 表示人名开始, \texttt{I-PER} 表示人名内部, \texttt{O} 表示实体外部.

深度学习方法的优势是特征学习能力强, 能利用大规模预训练语言模型提供的上下文语义信息.
但它同样依赖高质量标注数据, 在领域迁移、长尾实体和新实体上仍可能出错.

\section{词性标注}


\subsection{基本概念}


词性(part-of-speech, POS)指词汇的基本语法属性.
词性标注(part-of-speech tagging)是根据一个词在特定句子中的上下文, 为该词标注正确词性.

例如“研究”既可以作动词, 也可以作名词.
单独看这个词很难决定词性, 必须结合上下文:

\begin{Verbatim}[breaklines=true,breakanywhere=true,fontsize=\small]
我们 / 研究 / 自然语言处理
这项 / 研究 / 很重要
\end{Verbatim}

第一句中的“研究”更像动词, 第二句中的“研究”更像名词.
因此, 词性标注本质上也是一个上下文消歧问题.

\subsection{作用}


词性标注的作用包括:

\begin{itemize}
\item 为句法分析提供基础.
\item 为上层应用提供语法线索.
\item 在规则系统中提供等价类.
\end{itemize}

例如许多动宾短语都可抽象为 \texttt{V + N}:

\begin{Verbatim}[breaklines=true,breakanywhere=true,fontsize=\small]
吃 / 晚饭
打 / 篮球
\end{Verbatim}

如果系统知道某个词是动词、某个词是名词, 就可以用更一般的规则描述一批相似结构, 而不必逐词列举.

\subsection{主要困难}


汉语词性标注的困难主要包括:

\begin{itemize}
\item 缺乏形态变化.
\item 常用词兼类现象严重.
\item 词性划分标准尚有分歧.
\end{itemize}

印欧语系中, 数、格、时态等往往会引起词形变化, 因而能为词性判断提供线索.
汉语词形通常不随这些语法范畴变化, 所以不能简单从形态变化判断词性.

同时, 汉语中兼类词很多.
课件指出, 汉语中相当一部分词具有多个词性, 而且越常用的词越容易兼类.
这使得词性标注必须依赖上下文模型.

此外, 不同语料库和标注体系对词性类别划分并不完全一致.
例如:

\begin{itemize}
\item LDC 的 CTB 语料库将汉语词性划分为 33 类.
\item 北大语料库将汉语词性划分为 26 类一级词类.
\item 其他研究也可能采用 25 类或其他粒度的划分.
\end{itemize}

因此, 词性标注结果必须和具体 tagset 绑定理解.

\subsection{词性标记集}


词性标记集(POS tagset)是标注时可选择的词性类别范围.
例如:

\begin{Verbatim}[breaklines=true,breakanywhere=true,fontsize=\small]
{ 动词 V, 名词 N, 形容词 A, ... }
\end{Verbatim}

词性标记集通常来源于语言学分类, 但 NLP 任务往往会进一步细化.
例如英文中可以把动词进一步区分为第三人称单数动词、过去式动词、现在分词、过去分词等.

课件提到的常见英文标注集包括:

\begin{itemize}
\item Brown 标注集, 共 86 个标记.
\item Penn Treebank 标注集, 共 45 个标记.
\item CLAWS C5 标注集, 共 62 个标记.
\end{itemize}

常见中文标注集包括:

\begin{itemize}
\item LDC CTB 标注集, 共 33 个标记.
\item 北大标注集, 包含 26 个一级词类标记和 106 个全部标记.
\item 中科院 ICTCLAS 词性标注集, 共 39 个标记.
\end{itemize}

不同标注集的粒度不同.
粒度越细, 表达能力越强, 但标注难度和模型学习难度也越高.

\subsection{基于规则的方法}


词性标注中的规则方法可分为两类:

\begin{itemize}
\item 基于人工设计规则的方法.
\item 基于自动提取规则的方法.
\end{itemize}

人工规则方法依赖语言学专家手写规则.
例如早期 TAGGIT 系统就使用大量上下文规则来消除兼类词歧义.

自动提取规则的代表是基于转换规则的错误驱动学习方法.
其基本流程是:

\begin{enumerate}
\item[1.] 使用初始标注器得到一份自动标注结果.
\item[2.] 将自动标注结果与人工标准答案比较.
\item[3.] 根据预先设计的规则模板提取候选转换规则.
\item[4.] 选择使错误数减少最多的规则加入规则库.
\item[5.] 用规则库修正当前标注结果.
\item[6.] 若还能找到有效规则, 则继续迭代; 否则终止.
\end{enumerate}

这种方法保留了规则的可解释性, 同时又能从数据中自动学习具体规则.
它的核心思想是: 每一步都选择最能减少当前错误的转换规则.

\subsection{基于统计的方法}


基于统计的方法把词性标注看作序列建模问题.
给定词序列:
\[
W = (w_1, w_2, \dots, w_T)
\]

要求找到最可能的词性序列:
\[
T^* = \arg\max_{T} P(T | W)
\]

使用贝叶斯公式可转化为:
\[
T^* = \arg\max_{T} P(W | T) P(T)
\]

在 HMM 词性标注中:

\begin{itemize}
\item 隐状态是词性标签.
\item 观测是词语.
\item 状态转移概率表示相邻词性之间的搭配规律.
\item 发射概率表示某个词性生成某个词的概率.
\end{itemize}

例如:
\[
P(T) \approx \prod_{i=1}^T P(t_i | t_{i-1})
\]
\[
P(W | T) \approx \prod_{i=1}^T P(w_i | t_i)
\]

于是目标变为:
\[
T^*
  = \arg\max_{t_1, \dots, t_T}
    \prod_{i=1}^T P(t_i | t_{i-1}) P(w_i | t_i)
\]

这个最优标签序列可以用 Viterbi 算法求解.
这和前面 HMM 章节中的解码问题是一致的: 给定观测序列, 求最可能的隐藏状态序列.

HMM 词性标注器通常要解决三个问题:

\begin{itemize}
\item 如何设计模型框架: 明确状态、观测、转移概率和发射概率.
\item 如何训练模型: 从已标注语料中估计转移概率和发射概率.
\item 如何进行标注: 用 Viterbi 算法寻找最优词性序列.
\end{itemize}

早期 CLAWS 系统就是统计词性标注的重要代表.
课件提到, 它使用词类共现概率矩阵来消除兼类词歧义, 自动标注正确率可达到较高水平.
后续统计方法逐步发展到以 HMM、CRF 和神经网络为代表的序列标注模型.

\section{词法分析的当前技术水平}


现代汉语词法分析通常不再把分词、NER、词性标注完全割裂开来.
它们都可以被抽象为序列标注任务:

\begin{itemize}
\item 分词: 为每个字标注词位.
\item NER: 为每个字或词标注实体边界和实体类型.
\item POS: 为每个词标注语法类别.
\end{itemize}

这种统一视角的好处是可以复用 HMM、CRF、BiLSTM、BERT 等序列建模技术.
在工程系统中, 也常把多个任务做成级联流程或多任务模型.

整体来看:

\begin{itemize}
\item 在新闻、百科等规范文本上, 分词和词性标注已经比较成熟.
\item 在人名、地名、机构名等常见实体上, NER 也能取得较好效果.
\item 在跨领域、口语化、网络化和新实体密集场景中, 仍然存在明显挑战.
\end{itemize}

词法分析的关键不只是选模型, 还包括:

\begin{itemize}
\item 标注规范是否清晰.
\item 词表和语料是否覆盖目标领域.
\item 是否能处理新词和专有名词.
\item 是否能和下游任务目标保持一致.
\end{itemize}

因此, 词法分析应被理解为 NLP 系统的基础层.
它的输出会直接影响后续模型能够看到怎样的词、实体和语法线索.
